\section{The selected one-point Core theorem on the actual patched objects}
\label{sec:selected-core-v040}

This chapter completes the join left explicit in CRW20. It uses the
actual universal rings and maps constructed in PRG, the actual
patched modules and support completions constructed in PMS, the
original limit of LIM18, and the actual full-block cotangent argument
of CRW. Its final scope is the selected one-point branch
\(h=0\), \(a=n=4\), \(k.g=0\), with
\(D=4[L^+:\mathbf Q]\) and \(D_0=D+4q\).
All earlier point data, both rank-four Links, the fixed complex
embedding, prime \(3\), and the original ideal \(J\) are retained.
The generic quotient-pair and point-lifting lemmas are proved at
their wider original scopes where stated.

% BEGIN SCI1
\paragraph{SCI1. The exact selected contract.}
\label{claim:SCI1}
Keep the original data of PRG and assume the antecedents of
\texttt{Core}: \(\ell\ne2\), \(T.\mathrm{ProL}\),
\(\mathrm{Adapted}(T,r)\), the actual
\(A:\mathrm{TWAut}(u,J,\mathrm{Unit},D,T)\), and
\[
 \mathrm{GenBy}\bigl(O,\mathrm{pairVals}(u)\cup
       \{A.b.\lambda(X_{\mathrm{unit}})\}\bigr).            \tag{SCI.1}
\]
Take its one original point \(x:\mathrm{Pt}(T,u,D_0)\), with
\(a=n=4\), \(k.g=0\), \(D_0=D+4q\), and the normalization
\[
 O\longrightarrow x.K\text{ is onto},\qquad
 x.f(A.b.\lambda(X_{\mathrm{unit}}))=x.t.                  \tag{SCI.2}
\]
The point contains its actual full determinant representation,
\(\mathrm{Kap}\), auxiliary eigenvalues, and the whole B1/B2
statements. In particular B1 has its one constant before every
truncation and the original eventual level sets.
Here \(k=x.k\); for readability \(x.K,x.t,x.\rho\) denote the
field, parameter and representation of the underlying \(x.x\),
while \(x.\mathrm{al}\) is the original point's eigenvalue list.

Write \(p=\ker x.f\). The exact requested one-point
\(\mathrm{ChainConcl}\) says: if there is a prime
\(q_1\subseteq p\) with \(J\subseteq q_1\) and
\(\mathrm{DiscSt}(u.D,q_1)\), then every prime
\(q\subseteq p\) with \(\mathrm{StOrTh}(u,q)\) contains
\(J\). The intermediate-prime function has domain
\(\mathrm{Fin}(0)\), so there is no intermediate prime or
non-semisimplicity condition to construct. The final argument
actually proves \(J\subseteq q\) for every prime below \(p\)
once the stated starting prime is supplied. The original stable
and stable-or-thick hypotheses are retained in the theorem.
% END SCI1

% BEGIN SCI2
\paragraph{SCI2. A checked multiplicity-free Cayley--Hamilton structure theorem.}
\label{claim:SCI2}
Let \(C\in\mathcal C_O\), let \(H\) be a group, and let
\(\mathcal D\) be a positive-degree full determinant on
\(C[H]\) whose residual law is the determinant of distinct
characters \(\bar\chi_1,\ldots,\bar\chi_n:H\to k^\times\).
Let
\[
 E=C[H]/\mathrm{CH}(\mathcal D),                           \tag{SCI.3}
\]
where the ideal contains all coefficients of all generic
Cayley--Hamilton identities. The law descends to \(E\): the
checked kernel and Cayley--Hamilton statements in PQL7 give
\(\mathrm{CH}(\mathcal D)\subseteq\ker\mathcal D\).
The descended law is Cayley--Hamilton by construction.

We apply Chenevier's Theorem 2.22(ii), pp.~34--35 of the version
cited in PRG5. Its base is henselian because \(C\) is complete
Noetherian local. This henselian assertion follows by successively
correcting coprime monic factors modulo powers of the maximal
ideal: a Bezout identity for the two residual factors solves
each linear correction, and completeness takes the limit.
The residual law is split by the displayed diagonal representation
and is multiplicity free. Distinct characters remain distinct on
the residual Cayley--Hamilton quotient. Indeed each character
kills the Cayley--Hamilton ideal, and distinct character algebra
maps have distinct maximal kernels; the Chinese remainder theorem
gives the product of their copies of \(k\). All their dimensions
are one. The theorem therefore supplies orthogonal idempotents
\[
 e_1+\cdots+e_n=1,\qquad e_iEe_i=Ce_i,\qquad
 \operatorname{tr}_{\mathcal D}(c e_i)=c.                 \tag{SCI.4}
\]
The corner identities hold over \(C\), and tensoring them
preserves the needed identities over every scalar algebra.
No finite generation of \(E\), hypothesis \(\ell>n\), or
factorial inversion occurs in this application. It is the
multiplicity-free part of the published theorem, not its
absolutely irreducible part. The elementary consequences needed
for descent are proved in SCI3 rather than hidden inside a
trace-determines-determinant assertion.
% END SCI2

% BEGIN SCI3
\paragraph{SCI3. Trace-kernel equals full determinant kernel in this setting.}
\label{claim:SCI3}
For the determinant of SCI2,
\[
 z\in\ker\mathcal D
   \quad\Longleftrightarrow\quad
 \operatorname{tr}_{\mathcal D}(y z)=0
          \text{ for every }y\in C[H].                    \tag{SCI.5}
\]
It suffices on the right to test all single group elements,
since the trace is \(C\)-linear. Linearity and the cyclic
identity for the trace follow by comparing the degree-one term
in multiplicativity over a dual-number ring; those comparisons
involve no division. The forward implication of (SCI.5) follows
by taking that term in the full kernel criterion of PQL7.

For the converse pass to (SCI.3), and denote the image of \(z\)
again by \(z\). For every \(i\) and \(y\in E\), (SCI.4)
and the cyclic trace identity give
\[
 e_i y z e_i=\operatorname{tr}_{\mathcal D}(e_i y z)e_i=0.
 \tag{SCI.6}
\]
We use the integral identity
\(\mathcal D(1+ab)=\mathcal D(1+ba)\), valid after every
scalar extension. Here is a proof. In one extra polynomial
variable put \(u(t)=1+t(a-1)\). From
\((1+u b)u=u(1+b u)\) and multiplicativity, cancellation of
\(\mathcal D(u)\) gives
\(\mathcal D(1+ub)=\mathcal D(1+bu)\): this cancellation is
valid because \(\mathcal D(u)\) has constant coefficient \(1\)
and is thus a non-zero-divisor in the polynomial ring. Evaluating
at \(t=1\) proves the identity.

Fix a scalar algebra \(S\), extend the idempotents to \(E_S\),
and let \(w=h(ze_i)\) for arbitrary \(h\in E_S\).
Then \(w=we_i\). By (SCI.6), extended linearly to tensors,
\(e_iw=0\). Thus
\(\mathcal D_S(1+w)=\mathcal D_S(1+we_i)
 =\mathcal D_S(1+e_iw)=1\).
The full kernel criterion, quantified over every \(S,h\),
implies \(ze_i\in\ker\mathcal D\). The kernel is a two-sided
ideal and \(z=\sum_i ze_i\), so \(z\) belongs to it.
Pullback through (SCI.3) proves (SCI.5). This is the exact
\texttt{oo\_kills} implication needed below, with its
multiplicity-free hypothesis checked rather than omitted.
% END SCI3

% BEGIN SCI4
\paragraph{SCI4. The full determinant descends to the actual limiting global group.}
\label{claim:SCI4}
Choose the actual \(\Lambda\) by LIM18 and \(W,m\) by PRG13
and PRG24. Let \(I=\mathrm{ii}(W)\), the ideal generated by
all six original families \(\mathrm{Rel}\), and set
\(C=W.R/I\). PMS30, the proof of \texttt{piKer} gives
\(I\subseteq\ker(m.\pi)\); hence \(m.\pi\) factors as
\[
 \bar\pi:C\longrightarrow S_0=\mathrm{Spo}(A).
 \tag{SCI.7}
\]
Since \(S_0\ne0\), \(I\) is proper and PRG3 makes
\(C\) an object of \(\mathcal C_O\).

Push \(W.t.D\) to \(C\), obtaining \(\mathcal D_C\) on
\(\widetilde G=\mathrm{Ph0}(\Lambda)\). Its residual characters are
the distinct ones checked in PRG7. For every
\(k\in\ker\pi_{o,0}\), the original relation
\(\mathrm{Rel.ker}(h,k)\) gives
\[
 \operatorname{tr}_{\mathcal D_C}([h]([k]-1))=0
       \quad(h\in\widetilde G).
\]
SCI3 therefore says that the \emph{full} determinant kills
\([k]-1\). The kernel of the surjective group-algebra map
\(C[\widetilde G]\to C[\Lambda.G.G_0]\) is generated by these elements:
choose representatives of the group fibers and subtract their
coefficients, or perform that calculation on each finite-support
sum. The same statement holds over every scalar algebra.
PQL7's factorization criterion consequently produces a unique
full law \(D_C\) on \(\Lambda.G.G_0\) with
\[
 D_C\circ\pi_{o,0}=\mathcal D_C.                          \tag{SCI.8}
\]
It is continuous. Each of its characteristic coefficients modulo
\(\mathfrak m_C^a\) becomes locally constant after composition
with \(\pi_{o,0}\). This map is a continuous surjection of
profinite groups and hence a quotient map; the coefficient is
therefore locally constant itself. This proves the original
\texttt{ooDet} on the actual quotient ring.
% END SCI4

% BEGIN SCI5
\paragraph{SCI5. The complete original \(\mathrm{OoPair}\) and its specialization.}
\label{claim:SCI5}
Pull the residual characters of \(r\) back by the actual
surjection \(\Lambda.h.p_0\), and retain the identical local
residual characters and \(\theta\). This gives the original
\(r'\) and \(\mathrm{Compat}(\Lambda.h,r',r)\): the characters
remain distinct and continuous, their diagonal law has the
required polarization because the morphism commutes with the
conjugation and cyclotomic character, and the local groups and
\(\mu_v\) are unchanged. Take \(D_C\) from SCI4 and the
pushforwards of \(W.t.\psi\) to \(C\).

Their continuity, residual and \(\mu_v\) fields follow from
PRG2 and SCI4. The remaining \(\mathrm{Pair}\) fields follow
from the actual relations, as follows. To prove polarization at
a group element, lift it to \(\widetilde G\), conjugate the lift by the
given \(\Lambda.\mathrm{ct}\), and use \(\mathrm{Rel.pol}\).
Its image is the original conjugate since
\(\pi_o(\Lambda.\mathrm{ct})=\Lambda.G.c\); the cyclotomic
factor maps exactly by the \(O\)-algebra structure. For
ordinarity lift all the prescribed \(\Lambda.G.\phi_v(g_i)\)
and the testing group element through \(\pi_{o,0}\).
The image of \(\mathrm{Rel.ord}\) is exactly the required
ordered-product trace relation. For inertia write the given
element of \(r'.I_p(w)\) as \(\Lambda.G.\phi^u_w(g)\)
with \(g\in L.J_u(w)\), lift it and the testing element,
and use \(\mathrm{Rel.ine}\). Membership in \(r'.I_p\)
includes this \(I_u\) representation. All these lifts exist
by the same surjection; their choices disappear in the quotient.

This constructs every field of the original \(o:\mathrm{OoPair}(W)\).
Furthermore
\[
 o.p.D\otimes_{\bar\pi} S_0
       =(u.D\circ\Lambda.h.p_0)\otimes_B S_0,\qquad
 \bar\pi(o.p.\psi)=\iota_B(u.\psi).                       \tag{SCI.9}
\]
The character equality is the corresponding PRG23 field. For
the determinant, pull both sides back along \(\pi_{o,0}\),
use (SCI.8) and the full \(m.\pi_{\mathrm{maps}}\) equality
of PRG23, and cancel this pullback. Cancellation is valid for
full laws because the group-algebra map is onto after every
scalar extension. Thus (SCI.9) proves the actual
\texttt{ooMaps}, not merely equality of traces.
% END SCI5

% BEGIN SCI6
\paragraph{SCI6. The original quotient is generated by its actual point data.}
\label{claim:SCI6}
The ring \(W.R\) satisfies \(\mathrm{GenBy}\) for the union
of all its characteristic coefficients on group elements,
all ordered-character values, all \(U_{j,i},y_{j,i}\), and
\(W.t.s\). To prove it without changing \(W\), use the
integral word identity in JSP9: equality of all characteristic
coefficients on group elements determines a full determinant
law, including after scalar extension and in characteristic
\(3\). Two \(O\)-algebra maps agreeing on the displayed set
therefore carry all five tuple entries to the same targets;
\(W.\mathrm{univ}\) makes them equal.

For clarity, the generation consequence of this uniqueness is
elementary. Let \(R=W.R\), let \(R_0\) be the ordinary
\(O\)-subalgebra generated by the displayed values, and put
\(I_0=(R_0\cap\mathfrak m_R)R\). If
\(\mathfrak m_R/(I_0+\mathfrak m_R^2)\ne0\), a nonzero
\(k\)-linear functional on it defines an \(O\)-derivation
\(\delta:R\to k\). Subtract an \(O\)-lift of the residue
before applying the functional; the choice is immaterial because
\(\mathfrak m_OR\subseteq I_0\). The Leibniz identity follows
by discarding products in \(\mathfrak m_R^2\). The two maps
\(b\mapsto\bar b\) and
\(b\mapsto\bar b+\delta(b)\varepsilon\) to
\(k[\varepsilon]/(\varepsilon^2)\) agree on \(R_0\) but
differ, a contradiction. Thus the displayed quotient is zero.
Choose finitely many elements of \(R_0\cap\mathfrak m_R\)
spanning \(\mathfrak m_R/\mathfrak m_R^2\); monomials in
them span each graded piece. Correct a scalar approximation
one degree at a time to approximate any element modulo any
\(\mathfrak m_R^N\) by \(R_0\).

Passing to the surjective quotient \(R\to C\) preserves
these approximations. By (SCI.8), the determinant coefficients
become those of \(o.p.D\) at the images of the same group
elements; all character values are those of \(o.p.\psi\).
The other values become the actual quotient \(U,y,s\).
Consequently this is precisely the \(\mathrm{gen}\) field
required by \(\mathrm{OoPt}\). The integral word identity is
used only to identify full laws; no trace-only uniqueness
assumption is introduced.
% END SCI6

% BEGIN SCI7
\paragraph{SCI7. One splitting prime, with both endpoint primes below it.}
\label{claim:SCI7}
Since \(S_0\) is finite free of positive rank over \(B\)
by PRG18, \(B\to S_0\) is faithfully flat. There is a prime
\(P\subset S_0\) with \(P\cap B=p\), and for every prime
\(q'\subseteq p\) there is a prime \(Q'\subseteq P\)
with \(Q'\cap B=q'\). These are exact contractions.
We use the Stacks Project, Lemmas 10.39.16, 10.39.17 and
10.39.19, Tags 00HQ, 00HR and 00HS:
\url{https://stacks.math.columbia.edu/tag/00HS}.
The required flatness is the finite freeness just proved;
faithfulness follows because tensoring with a nonzero finite
free module preserves a nonzero module. To see the specified
going-down application, localize at \(p,P\). The map
\(B_p\to(S_0)_P\) is flat and local, hence faithfully flat;
its spectrum surjects onto \(\operatorname{Spec}B_p\).
Pull a prime over \(q'B_p\) back to \(S_0\).

Choose \(P\) once, before the two endpoint primes are lifted.
Both \(q_1\) and the target \(q\) of SCI1 can therefore be
lifted below this same \(P\). The later completion and the
later point will use this very \(P\); choosing unrelated
splitting primes for the two endpoints would not suffice.
% END SCI7

% BEGIN SCI8
\paragraph{SCI8. A lift of the original point with exactly the prescribed kernel.}
\label{claim:SCI8}
For every prime \(P\subset S_0\) above \(p=\ker x.f\),
there is an \(O\)-map
\[
 f':S_0\longrightarrow A_x=x.K[[t]],\qquad
 f'|_B=x.f,\qquad\ker f'=P.                               \tag{SCI.10}
\]
Here is a direct proof of the source's \texttt{lift\_pt}.
Let \(B_0=B/p\) and \(E_0=S_0/P\), both domains.
The latter is finite integral over the former and is generated
by the images of the actual roots of PRG18. Put
\(F_0=\operatorname{Frac}B_0\),
\(E=\operatorname{Frac}E_0\), and
\(F_x=\operatorname{Frac}A_x\); the injective map induced by
\(x.f\) embeds \(F_0\) into \(F_x\). The extension
\(E/F_0\) is finite. The nonzero finite \(F_x\)-algebra
\(E\otimes_{F_0}F_x\) has a maximal ideal. Its residue is
a field \(M\) containing both \(E\) and \(F_x\), since
the maps from fields to a nonzero ring are injective.

Each generating root satisfies the original characteristic
polynomial, whose image in \(A_x[X]\) is
\[
 \operatorname{charpoly}(x.\rho(T.\mathrm{sg}_j))
       =\prod_i(X-x.\mathrm{al}_{j,i})                    \tag{SCI.11}
\]
by the full point map and \(x.\mathrm{eig}\). In the field
\(M\), that root must equal one of these elements of \(F_x\).
Since these roots generate \(E\), the residue algebra \(M\)
is generated over \(F_x\) by elements already in \(F_x\).
Thus \(M=F_x\), giving an embedding \(E\to F_x\)
extending \(F_0\). Every element of \(E_0\) is integral
over \(B_0\), so its image is integral over \(A_x\).
The power-series DVR \(A_x\) is integrally closed: an element
of its fraction field of negative valuation cannot satisfy a
monic equation, by comparing the valuation of its highest power
with those of all lower terms. Hence the image of \(E_0\)
lies in \(A_x\). The resulting map (SCI.10) is injective
after quotienting by \(P\), which proves its exact kernel.
It is an \(O\)-map because it extends \(x.f\) on \(B\).
No new coefficient field, point kernel, or eigenvalue extension
has been chosen.
% END SCI8

% BEGIN SCI9
\paragraph{SCI9. The actual roots avoid the patching denominator.}
\label{claim:SCI9}
Put \(\alpha_{j,i}=f'(m.\pi(W.t.U_{j,i}))\).
For each \(j\), PRG23 and the ordered splitting relations give
\[
 \prod_i(X-\alpha_{j,i})
   =\operatorname{charpoly}(x.\rho(T.\mathrm{sg}_j))
   =\prod_i(X-x.\mathrm{al}_{j,i}).                        \tag{SCI.12}
\]
The right-hand roots are distinct by \(x.\mathrm{eig}\).
Over the domain \(A_x\), equality of products of linear factors
implies equality of their root multisets: evaluate at one root,
find the equal root on the other side, cancel its monic factor,
and proceed by induction on the number of factors. Thus the
\(\alpha_{j,i}\) are distinct as well. The original
\(f_o=\mathrm{fEl}(1,W.t.U)\) consequently satisfies
\[
 f'(m.\pi(f_o))
   =\prod_j\prod_i\prod_{i'\ne i}
          (\alpha_{j,i}-\alpha_{j,i'})^2\ne0.              \tag{SCI.13}
\]
An empty product is one, so this also covers \(q=0\).
Hence \(f_o\notin(m.\pi)^{-1}(P)\), exactly the avoidance
condition in S5. The point's original eigenvalue list and its
B1/B2 witnesses have not been relabeled: (SCI.12) only identifies
the possibly permuted roots required by the separate
\(\mathrm{OoPt}\) record.
% END SCI9

% BEGIN SCI10
\paragraph{SCI10. Every field of the actual lifted ordinary point.}
\label{claim:SCI10}
There is an original \(\mathcal P:\mathrm{OoPt}(\Lambda.h,x.x,o.p)\)
with
\[
 \mathcal P.f=f'\circ\bar\pi,\qquad
 \mathcal P.\mathrm{ker}=\mathrm{ooKer}(o).                \tag{SCI.14}
\]
We check every field. The finite topological generating set is
the image of \(\Lambda.\mathrm{fg}\) under \(\pi_o\): a
continuous surjection maps the closure of the generated subgroup
onto a dense subgroup of the target. The equality
\(r'.I_p(w)=\Lambda.G.\mathrm{frame}.I_u(w)\) follows from
adaptedness. For an element of the latter group, write it using
the same local \(g\in L.J_u(w)\); the morphism sends it to
an element of the original \(r.I_p(w)=G_0.I_u(w)\), and the
pulled-back residual characters are all one. The reverse
inclusion is definitional. Surjectivity of \(O\to x.K\)
is precisely (SCI.2).

Use the literal \(q=T.q\), \(\mathrm{ph}=\Lambda.\mathrm{ph}\),
\(\mathrm{ta}=\Lambda.\mathrm{ta}\). Their commutativity
and \(\Lambda.h.p_0(\mathrm{ta})=1\) are fields of the same
limit, supplied by LIM18. Take \(U,y,s\) to be the images
in \(C\) of the actual \(W.t.U,W.t.y,W.t.s\).
The q1 and q2 trace identities follow by lifting their testing
element, \(\mathrm{ph}\), and \(\mathrm{ta}\) through
\(\pi_{o,0}\), and then killing exactly \(\mathrm{Rel.q1}\)
and \(\mathrm{Rel.q2}\). Each ordered product maps factor by
factor, so this preserves its specified order. The generation
field is SCI6.

Take the map to be \(f'\bar\pi\). Its full determinant and
all ordered-character identities follow from (SCI.9), (SCI.10),
and the original \(x.\mathrm{maps}\). Its value on \(s\)
is \(t\) by PRG23 and (SCI.2), and its value on every \(y\)
is zero by PRG23. Injectivity of each root list is SCI9.
The eigenvalue polynomial is (SCI.12), with
\(\Lambda.h.p_0(\Lambda.\mathrm{ph}_j)=T.\mathrm{sg}_j\)
from the same limit. These are all fields of \texttt{OoPt}.
Finally its closed relation set is defined by this very pair,
\(\mathrm{ph},\mathrm{ta},U,y\); it is literally the
\(\mathrm{ooKer}(o)\) in (SCI.14), with no ideal comparison
or replacement. This proves the original \texttt{ooPt}
construction on the chosen prime.
% END SCI10

% BEGIN SCI11
\paragraph{SCI11. Match the actual S5 completion to the actual ordinary point.}
\label{claim:SCI11}
Set
\[
 \mathfrak p_*=(m.\pi)^{-1}(P)\subset W.R,\qquad
 \bar{\mathfrak p}_*=\mathfrak p_*/I\subset C.
\]
SCI4 shows \(I\subseteq\mathfrak p_*\); SCI9 gives
\(f_o\notin\mathfrak p_*\); and PRG23 gives
\(m.s(Y_{j,i})\in\mathfrak p_*\). These are exactly the
prime, relation, denominator and parameter guards of S5.
Moreover, by (SCI.10) and (SCI.14),
\[
 \ker\mathcal P.f=\bar\pi^{-1}(P)=\bar{\mathfrak p}_*.
 \tag{SCI.15}
\]
PMS29 applies its actual finite modules
\(M_N=(A.A_N).M\), the actual \(\rho_N\)-actions, the one
ultrafilter \(\Lambda.U\), and the actual augmentation fiber
\(X_0=S_0\otimes_B A.b.M\). Its freeness, paired quotient
maps and all six approximate-annihilator families have been
proved there from the original \(T,A\); they are not new
premises in this chapter. The normalized coefficient-prime
completion has dimension \(1+D\): its defining map is
\(x.f\circ A.b.\lambda\), since PRG23 and (SCI.10) give
\[
 m.\pi(m.s(Ca))\in P
       \quad\Longleftrightarrow\quad
 x.f(A.b.\lambda(a))=0\qquad(a\in L_0).                  \tag{SCI.16}
\]
PMS16, the coefficient-completion argument, uses exactly the
surjectivity and \(s\mapsto t\) of (SCI.2).

Thus PMS15 and PMS29 give a completion \(C_*\) of the actual
\(C_{\bar{\mathfrak p}_*}\), a support ideal \(S_*\), and
the precise fiber equivalence
\[
 S_*\subseteq Q
 \quad\Longleftrightarrow\quad
 \operatorname{Ann}_{W.R}(X_0)
       \subseteq\operatorname{contraction}_{W.R}(Q)        \tag{SCI.17}
\]
for every prime \(Q\) containing the parameter images.
Every minimal prime over \(S_*\) has quotient dimension at
least
\[
 (1+D)+\#(\mathrm{Fin}(q)\times\mathrm{Fin}(n))
        =1+D+nq=1+D_0.                                   \tag{SCI.18}
\]
By (SCI.15), this is precisely the completion and ordinary
point to which CRW applies. LIM19, equivalently CRW12--CRW15
on the actual limit, supplies K1 from the same point's B1.
No unrelated completion, Selmer group or parameter count is
introduced in this match.
% END SCI11

% BEGIN SCI12
\paragraph{SCI12. The selected normalized Core theorem is complete.}
\label{claim:SCI12}
Under the original hypotheses in SCI1, with the actual constructions
of SCI4--SCI11,
the one-point \(\mathrm{ChainConcl}(u,J,p)\) holds.
Take its original starting prime \(q_1\) and any target prime
\(q\subseteq p\). By SCI7 choose \(Q_1,Q\subseteq P\)
in \(S_0\) contracting respectively to \(q_1,q\), with
the single \(P\) already used in the lifted point.
The original \texttt{piSurj}, proved in PMS31 using (SCI.1),
and the actual fiber-support calculation PMS32 give
\[
 \operatorname{Ann}_{W.R}(X_0)\subseteq(m.\pi)^{-1}(Q')
       \quad\Longleftrightarrow\quad J\subseteq Q'\cap B
       \quad(Q'\text{ prime in }S_0).                    \tag{SCI.19}
\]
This is the exact original ideal \(J\), from
\(A.b.\mathrm{supp}\). In particular the starting prime
contains the annihilator after contraction.

Contract \(Q_1\) through \(\bar\pi\), localize at
\(\bar{\mathfrak p}_*\), and lift it to a prime of \(C_*\).
This lift exists by faithful flatness of completion of a
Noetherian local ring, the checked application already in CRW19.
All parameter images vanish by PRG23. Hence (SCI.17) makes
this prime contain \(S_*\). Choose a minimal prime \(Z\)
over \(S_*\) below it. By (SCI.18),
\(\dim(C_*/Z)\ge1+D_0\). The original full-block cotangent
presentation of CRW2--CRW18, applied through SCI10--SCI11,
is a surjection from a regular power-series domain of dimension
\(1+D_0\). CRW18 therefore forces \(Z=0\).
Consequently \(S_*=0\).

Now contract the target \(Q\) through \(\bar\pi\), and
lift it to a prime of the same \(C_*\) by the same faithful
flatness. It contains all the parameters, and trivially contains
\(S_*=0\). Equation (SCI.17) then gives its contracted
annihilator containment; (SCI.19) gives
\(J\subseteq Q\cap B=q\). This is the required conclusion.

The stable condition at \(q_1\) and the stable-or-thick
condition at \(q\) remain among the original antecedents.
The full-block dimension argument made them unnecessary within
this selected proof. Accordingly no reconstruction of the
general \texttt{NonssT} comparison or of a multi-point move
is needed for this conclusion. This is a complete selected
\(\mathrm{Core}\) producer, not an application with assumed
\(\mathrm{RRObj}\), \(\mathrm{PMod}\), \(\mathrm{OoPt}\),
or assumed support data. Their constructions have all been
supplied on the original objects.
% END SCI12

% BEGIN SCI13
\paragraph{SCI13. Normalize and descend: the selected Raw endpoint is complete.}
\label{claim:SCI13}
Start with the original selected \(h=0,a=n=4,k.g=0\)
\(\mathrm{Raw}\) data, whose automorphic packet has the
\(\mathrm{Empty}\) parameter type. PQL constructs the actual
\(\mathrm{TWQuot}\), descends its full pair, point, Kap,
eigenvalues, B1/B2 and given \(\mathrm{TWAut}\), and proves
\(\mathrm{ProL}\) for that quotient system. Its original
\(B,J,K,f\) and its finite modules are unchanged; all
constant-before-truncation quantifiers are preserved.

Apply the actual coefficient and Unit normalization of BCN24
to this descended one-point datum. It constructs its finite
coefficient extension and the literal enlarged ring and module
packet, maps the new Unit variable to \(t\), and makes the
coefficient map onto the same point field. The values
\(h=0,a=n=4,k.g=0\) are unchanged. BCN24 invokes only the
restricted normalized Core conclusion proved in SCI12; it
does not require a theorem for all Kap or arbitrary chain length.
Its exact prime-contraction argument brings \(J\)-containment
back to the original ring. The pointwise matched descent in
PQL20--PQL21, equivalently the explicit composition in BCN25,
then gives the original selected \(\mathrm{Raw}\) conclusion.

Thus the limiting group, universal ring/maps, actual limiting
module, support/completion theorem, quotient pair and lifted
ordinary point, coefficient normalization/Unit, and pro-\(\ell\)
descent are all supplied for this selected route. The current
arithmetic obstruction is earlier: constructing the original
ordinary/control/type data, the four Hida endpoints and the
all-level same-\(J\) \(\mathrm{TWAut}\)/\(\mathrm{HasAut}\)
packet on both original rank-four Links with one common
coefficient threshold. This chapter consumes such a packet and
does not construct it. The whole compound Count-II obligation,
and therefore the complete human RH/Dirichlet-GRH proof, remain
open. The claim is the completed selected Raw producer with
its exact original arithmetic antecedents.
% END SCI13
